\section{Three Response Models and Their Tests}

The medium suggests mechanisms to investigate. The task and data determine which response model fits.

\subsection{Approach to a Usage Plateau}

For a defined digital service under sufficiently stable conditions, an illustrative model is
\begin{equation}
 Y_{\mathrm{bits}}(t)=Y_\infty(1-e^{-t/\tau}).
\end{equation}
The initial value is zero. The response approaches a fitted plateau $Y_\infty$ over timescale $\tau$. Demand, budgets, distribution, usefulness, and capacity can all affect the parameters.

Measure actual use or effective service, not merely generated volume. Compare the exponential approach with alternatives and re-estimate after material changes in the environment.

\subsection{A Capacity-Limited Physical Process}

For fixed facilities and a consistently defined input $I$, one possible model class satisfies
\begin{equation}
 0\leq Y_{\mathrm{atoms}}(I)\leq\Phi_{\max}.
\end{equation}
If diminishing marginal gains are additionally proposed over the observed range,
\begin{equation}
 \frac{\partial^2Y_{\mathrm{atoms}}}{\partial I^2}\leq0.
\end{equation}
The upper bound and concavity are separate assumptions. New capacity, learning, complementarities, or step changes can violate the latter and change the former.

Measure design, fabrication, construction, verification, and maintenance separately. AI can improve one stage without determining the rest. Complete-task measurements are needed before comparing its effect across industries.

\subsection{An Adaptive Response}

For one specified biological or adaptive process, a candidate observation is
\begin{equation}
 \frac{\mathrm d\langle Y\rangle}{\mathrm dP}\approx0,
 \qquad
 \frac{\mathrm d\,\operatorname{Var}(Y)}{\mathrm dP}>0,
 \quad P>P_c.
\end{equation}
Mean benefit reaches a plateau while dispersion increases beyond a fitted threshold. Other processes can improve, decline, or respond after a delay. A variance increase is not automatically an increase in harmful outcomes; the outcome definition must identify the risk.

Collect input intensity, duration, population, baseline state, repeated measurements, and follow-up. Distinguish dose from pace. Experiments or credible comparison designs must separate the intervention's effect from selection and changing conditions.

\subsection{A Common Evaluation Record}

The three examples use different independent variables: time, input, and intervention intensity. Their plotted slopes are not directly comparable.

To compare a shared technology across stages or media, fix a useful task, define comparable input and outcome measures, and record the resources, elapsed time, accepted results, failures, and maintenance required. Matched projects, staged deployments, and suitable experimental designs can then test particular mechanisms.

A multistage project's total output is generally not a linear sum of these three curves. Interactions and bottlenecks require a model of the actual chain.

The empirical program succeeds when it identifies which constraint changed, what happened downstream, and whether a competing explanation accounts for the same evidence more simply.
